\chapter{Self-Explanation and Learning}
\label{sec:p19-self-explanation-and-learning}

Explaining a worked step to oneself can expose a condition that the printed
solution leaves implicit. Explaining a sentence can connect it to an earlier
sentence, reveal a conflict with one's prior model, or produce a mistaken
inference. These are activities performed while learning, not guarantees
that the explanation is true or that the learner will later solve a new
problem. The empirical question is whether a specified explanation activity
changes a specified learning outcome, compared with a specified alternative.

This chapter's bounded primary corpus is Chi and colleagues' 1989 study of
physics examples, their 1994 study of elicited explanations of a biology
text, and Muldner, Burleson and Chi's 2014 diffusion-learning experiment
\cite{ChiEtAl1989SelfExplanations,ChiEtAl1994Eliciting,
MuldnerBurlesonChi2014Emergent}. The first principally compares learners
selected by subsequent performance; the later papers compare instructional
conditions and also make within-condition comparisons. These designs must
not be treated as interchangeable experiments. The evidence-record
distinctions in \cref{def:p19-empirical-question,def:p19-evidence-record}
apply. No pooled efficacy estimate or exhaustive literature review is
inferred from these three papers.

\section{What is being explained, and what counts as an explanation?}
\label{sec:p19-self-explanation-task}

\begin{hbdefinition}{def:p19-self-explanation-activities}{Explanation, paraphrase and monitoring}
In an instructional self-explanation task, a learner attempts to make sense
of material for themselves: why a solution step applies, what follows from
a sentence, or how a new assertion relates to information already available.
A \emph{paraphrase} restates supplied content without the additional
domain-relevant inference at issue. A \emph{monitoring statement} reports
understanding or difficulty; it does not itself supply a reason. Verbatim
rereading, planning to reread, and producing a substantive explanation are
also different observable activities.

These are operational distinctions, not one universal coding system.
The 1989 physics paper separates explanations, monitoring statements and
other ideas, including paraphrases. The 1994 paper's explanation counts
exclude paraphrases, monitoring and ordinary comprehension or bridging
inferences. The 2014 paper initially includes paraphrases in its broad
explanation category, then separately identifies explanations containing
additional inferences. Counts from those definitions cannot be pooled
as if they measured identical units
\cite[pp.~158--164]{ChiEtAl1989SelfExplanations}
\cite[pp.~454--457]{ChiEtAl1994Eliciting}
\cite[pp.~850--852]{MuldnerBurlesonChi2014Emergent}.
\end{hbdefinition}

An utterance's length is not its explanatory value. The physics coders
first counted lines and then grouped them into ideas; several lines could
develop one idea. The biology analysis compared coarse knowledge-inference
units with a finer proposition-level recoding. Repeatedly stating one
inference did not create several distinct coarse inferences. For diffusion,
the usual unit was the response to a prompt, and an inter-level code required
an explicit connection between visible behavior and molecular behavior.
The text mentioning both levels did not automatically make the learner's
response an inter-level explanation.

Nor is being coded as an inference equivalent to being correct. A learner
may add an unsupported causal claim. Content, relation to the instructional
material, correctness and grain size all belong in the coding record.
The biology paper reports two project-member coding passes with broadly
similar patterns, not a stated blinded reliability coefficient for every
analysis. The diffusion paper does not report coder blinding or an
interrater-reliability coefficient. Those omissions limit reproducibility
of the process measures; they do not license inventing a coefficient.

\begin{hbexample}{def:p19-self-explanation-incline}{A step, its reason and its limit}
Consider a worked mechanics example that chooses one coordinate direction
along a straight incline and the other perpendicular to it. A paraphrase
says that the axes have been placed along and across the slope. An
explanation adds why that choice is useful: under the example's stated
geometry, the normal contact force lies along the perpendicular direction,
so it need not be split into two components; weight still must be resolved.
The explanatory addition connects a representational choice to the forces
and the algebraic work it simplifies.

A second check asks whether the coordinate choice makes the body stationary.
It does not. Force balance requires the example's actual acceleration and
force assumptions, not merely a conveniently drawn pair of axes. For an
accelerating body, replacing the net-force equation with zero would be an
incorrect explanatory addition. To adapt the example, identify the body,
list the forces on that body, state its motion, and only then choose axes.

This is a handbook reconstruction of the kinds of conditions, consequences,
goals and meanings analyzed in the physics protocols, not a participant
quotation or a newly conducted learning trial
\cite[pp.~164--168]{ChiEtAl1989SelfExplanations}.
\end{hbexample}

\section{Spontaneous explanation: the physics study}
\label{sec:p19-self-explanation-physics}

Chi and colleagues recruited ten paid campus volunteers, five male and five
female. Eight were working toward bachelor's degrees and two had additional
postgraduate psychology training. None had taken college physics; all but
one had taken high-school physics. The researchers deliberately sought a
range of abilities rather than a representative random sample. Learning
took place in the laboratory over several weeks, with total study times
between eight and twenty-nine hours and no class lectures
\cite[pp.~151--154]{ChiEtAl1989SelfExplanations}.

Students studied the measurement, vector and one-dimensional-motion
background chapters of Halliday and Resnick, then the particle-dynamics
chapter. Background errors led to correction attempts, provision of worked
solutions, and, if needed, experimenter explanation. Before studying the
three target worked examples, students had to answer the chapter's
declarative and qualitative questions correctly. The examples concerned
strings, an inclined plane and a pulley. Students could explain what they
understood after reading example lines; their talk was recorded. Thus
``spontaneous'' refers to the substantive explanations students generated
within an invited talk-aloud procedure, not to completely unobserved
ordinary classroom study.

The subsequent task comprised twelve problems classified as isomorphic,
four for each example type, and seven end-of-chapter problems classified
as far transfer by the paper's isomorphism criteria. Changes included
surface features, gravity and force direction. Students gave protocols and
received no feedback during these problem-solving tasks. Arithmetic slips
were not counted as errors, partial solutions received proportional credit,
and problems had equal weights. The detailed example-use analysis was
limited to the first three isomorphic problem protocols, although all
nineteen problem scores determined the performance grouping
\cite[pp.~154--158, 172--175]{ChiEtAl1989SelfExplanations}.

One student who generated no protocols and solved no problems was excluded,
leaving nine. The analysis then omitted the middle scorer and compared the
four highest with the four lowest scorers. The original ``Good'' and ``Poor''
labels therefore describe selected outcome groups, not assigned treatments
or permanent learner types. Two judges conducted scoring and protocol
classification, with reported mean interrater reliabilities across analyses
ranging from 86 to 94 percent.

\begin{table}[htbp]
\centering
\small
\begin{tabular}{p{.47\linewidth}rr}
\hline
Reported measure & Higher scorers & Lower scorers\\
\hline
Students in selected comparison & 4 & 4\\
Overall problem success & 82\% & 46\%\\
Isomorphic problem success & 96\% & 62\%\\
Chapter problem success & 68\% & 30\%\\
Physics explanations per student per example & 15.3 & 2.8\\
Mean example-study time, minutes & 13.0 & 7.4\\
Initially articulated law components, out of 12 & 5.5 & 5.5\\
Additional law components gained or accessed & 3.0 & .25\\
\hline
\end{tabular}
\caption{The 1989 physics comparison. Group membership was selected using
the performance outcomes, so the success-rate contrast is not a randomized
explanation effect. Source: pp.~158--168, Tables~2--4
\cite{ChiEtAl1989SelfExplanations}.}
\label{tab:p19-self-explanation-physics}
\end{table}

The number of physics explanations was associated with subsequent
isomorphic success, with reported \(r=.81,p<.02\). The qualitative
analysis distinguishes explicating application conditions, drawing
consequences, identifying a goal and interpreting a quantitative expression.
Higher scorers more often linked these explanations to the principles
introduced in the prose. A component analysis found equal initially
articulated Newton-law knowledge, but different additional components
appearing while examples were studied. The authors explicitly retain the
alternative that example study accessed previously unarticulated knowledge,
rather than created all of it. Declarative recitation is an imperfect
measure of usable prior understanding
\cite[pp.~163--169]{ChiEtAl1989SelfExplanations}.

Monitoring also differed. Higher scorers produced more statements of
comprehension failure, often about conceptual choices rather than only
algebraic manipulation. Their questions could identify a particular missing
reason. Both groups used examples on most isomorphic tasks: the reported
global counts were nine versus ten of twelve problems. The finer reference
analysis found more extended rereading by lower scorers, but no differences
in the numbers of copying/mapping and comparing/checking episodes. The
total reference-episode contrast, 2.7 versus 6.7 per problem, was only
marginal by the paper's reported criterion, \(p<.10\), not a conventional
\(.05\) finding. Successful learners were not learners who never consulted
an example
\cite[pp.~170--175]{ChiEtAl1989SelfExplanations}.

This study supplies a detailed association and a process analysis. It has
no randomized no-explanation control, and the higher scorers spent longer
studying examples. Their similar amount of talk during problem solving
argues against a simple general-verbosity account, but does not independently
equalize engagement, time, prior procedural understanding or motivation.
The paper's far-transfer label concerns different problems in the same
physics chapter. It does not establish transfer to an unrelated domain,
durable improvement on a later unprompted test, or general reasoning ability.

\section{Elicitation rather than learner selection: the biology text}
\label{sec:p19-self-explanation-biology}

The 1994 study asked whether explanation could be elicited, rather than
merely associated with a selected learner's success. Its main sample was
twenty-four eighth-grade volunteers from one local inner-city public school:
fourteen prompted students, seven boys and seven girls, and ten unprompted
students, five boys and five girls. Eleven pilot students were tested
separately to refine materials and estimate timing. None of the main
participants had taken a biology course. The researchers deliberately sought
a range of California Achievement Test (CAT) scores and found lower-scoring
students harder to recruit. The method does not report random assignment
\cite[pp.~443--451]{ChiEtAl1994Eliciting}.

The passage was an edited circulatory-system unit from Towle's high-school
biology text, presented as 101 sentences, one sentence per sheet. Figures
and excursions into related topics were removed. The prompted condition
combined initial instructions, general sentence-by-sentence explanation
reminders, clarification questions where an utterance was vague, and
twenty-two specific component-function probes. Notes and diagrams were
permitted but not systematically elicited. The control read the passage
twice without these explanation or function prompts.

Twice-reading was intended to approximate the pilot group's explanation
time. Actual means were two hours five minutes for prompted students
(range one hour twenty-seven minutes to two hours fifty-three minutes)
and one hour six minutes for controls (range twenty-two minutes to two
hours forty-seven minutes). Thus
the implemented conditions did not equalize time. The contrast is an
instructional package, including specific probes and experimenter
interaction, not an isolated randomized change in one internal process.
Controls were not forbidden to explain covertly, and their covert
explanations could not be counted.

The pretest included definitions of twenty-three terms, a drawing of blood
flow and half the questions in each of four categories. The posttest
repeated the terms and drawing and included the full question sets,
including pretested questions. Students could consult the text or their
notes on the posttest. The study ran in three or four sessions, each one
to three hours long, at least a week apart. This supports assessment across
the implemented spacing; it is not an immediate-only test, but neither
is it an unassisted long-term retention experiment
\cite[pp.~450--452]{ChiEtAl1994Eliciting}.

\begin{hbdefinition}{def:p19-self-explanation-question-types}{Four tests with different inference demands}
The biology paper designed fourteen questions in each of four categories.
Verbatim questions requested explicitly stated material. Comprehension
inferences required integrating supplied lines or paragraphs. Knowledge
inferences required additional knowledge, such as inferring a component's
function and relating it to the system. A fourth category applied the
learned system to hypothetical health-related questions. These are
instructional questions, not evidence of clinical diagnosis or treatment
competence. Answers were scored against researchers' ideal-answer templates,
with a maximum of six points per question and partial credit.

New inferential questions are not the same test as literal recall. Repeated
pretest questions are not wholly novel transfer items. Open-book answers
need not measure unaided retrieval. All three distinctions matter when
interpreting a gain
\cite[pp.~447--452]{ChiEtAl1994Eliciting}.
\end{hbdefinition}

Both conditions improved. The reported gains on the percentage-correct
scale were 32 percentage points for prompted students and 22 for controls;
the group-by-test interaction was \(F(1,22)=5.1,p<.05\). For the harder
knowledge-inference and application categories, gains were 22.6 versus
12.5 points, with reported \(t(22)=2.64,p<.01\). The overall group main
effect was not significant, \(F(1,22)<1\); this does not contradict the
reported difference in changes from pretest to posttest. CAT means were
87 versus 83 percent, not significantly different. A regression retaining
CAT scores still reported a prompting advantage, but a non-significant
CAT term does not establish absence of prior-ability influence
\cite[pp.~452--454, Figures~2--3]{ChiEtAl1994Eliciting}.

Within the prompted condition, the authors selected four high and four low
explainers by their explanation counts. The two very lowest participants
were excluded from that comparison because one had an outlying CAT score
and the other lacked a CAT score. This selection was not a second randomized
experiment. Coarse counts averaged 87 versus 29 inferences, including
incorrect ones; correct counts averaged 61 versus 14. Overall gains of
38 versus 27 points gave \(t(6)=2.14,p<.07\), which did not meet a
\(.05\) criterion. Gains for the harder categories, 33 versus 17 points,
gave \(t(6)=3.02,p<.05\). Removing answers obtained by consulting the
text changed the overall comparison, but that is a different, post hoc
outcome definition
\cite[pp.~457--459]{ChiEtAl1994Eliciting}.

\begin{table}[htbp]
\centering
\small
\begin{tabular}{p{.38\linewidth}p{.26\linewidth}p{.27\linewidth}}
\hline
Comparison & Reported result & Interpretation boundary\\
\hline
Prompted versus twice-reading & 32 versus 22 point gain;
\(F(1,22)=5.1,p<.05\) & Unequal actual study time; assignment mechanism
not reported as random\\
High versus low explainers, overall & 38 versus 27 point gain;
\(p<.07\) & Selected count extremes; not a \(.05\) result\\
High versus low, hard questions & 33 versus 17 point gain;
\(t(6)=3.02,p<.05\) & Eight selected students, different outcome subset\\
Newly attaining double-loop-2 model & Prompted 8/14; controls 2/9 &
Control denominator excludes a student already correct at pretest\\
Selected high versus low, final model & 4/4 versus 1/4 double-loop-2 &
Descriptive selected subgroup comparison\\
\hline
\end{tabular}
\caption{Distinct 1994 contrasts, not five independent experiments.
Reported gains are changes in percentage-correct scores.
Source: pp.~452--459, 467--469 \cite{ChiEtAl1994Eliciting}.}
\label{tab:p19-self-explanation-biology}
\end{table}

The model analysis reconstructed the connections and directions of blood
flow from the term interviews and drawings, then classified models by
explicit necessary features. Six types ranged from no loop to the paper's
most complete double-loop model. It was not simply a correct-answer total:
the top classification required evidence of the connections, chambers,
separation and routes specified in the coding scheme. Eight of fourteen
prompted students newly reached that classification, versus two of nine
eligible controls. The control posttest table contains three of ten at that
level, but one had already reached it at pretest. Comparing 8/14 with 2/9
without explaining that eligibility rule would confuse attainment with
unconditional final status
\cite[pp.~464--469, Tables~4--5]{ChiEtAl1994Eliciting}.

The paper also retains boundaries. No high/low difference was reported for
articulating explicitly stated functions; their difference was in inducing
unstated functions. The five higher-CAT and five lower-CAT prompted
students gained about the same amount, but this small selected comparison
does not prove equal benefits at all ability levels. Diagram drawing and
verbosity differences were not significant, and drawing was not assigned
as a diagram-only treatment. The observation that some controls drew many
diagrams suggests constructive engagement, not a direct measurement of
their covert self-explanation. Moreover, the source describes fourteen
questions in each of four categories, while its text-reference analysis
calls the posttest fifty-four questions. Without the item records, those
different counts cannot be silently reconciled
\cite[pp.~458--463]{ChiEtAl1994Eliciting}.

\section{A boundary case: learning molecular and visible diffusion}
\label{sec:p19-self-explanation-diffusion}

The 2014 experiment extends the question to diffusion, where visible
stability need not imply microscopic inactivity. Forty-two college students
participated for psychology credit. The paper describes a pretest-stratified
random procedure for equalizing prior knowledge between conditions, with
twenty-three self-explaining and nineteen reading students. Sessions were
individual and video-recorded.
Both conditions read aloud and could interact with molecular-level and
visible-level simulations. The explanation condition received fifty-two
generic text prompts, including nine warm-up prompts, and eleven specific
simulation prompts; the reading condition read twice and was invited to
explore the simulations freely
\cite[pp.~849--850]{MuldnerBurlesonChi2014Emergent}.

The comparison therefore changes both generic explanation and directed
simulation questions. Asking about molecules at equilibrium can direct
attention to information that free exploration misses. The authors
acknowledge this difference; the fact that specific prompts were a minority
does not experimentally exclude their importance. Mean study times are not
reported, so twice-reading establishes a time-control intention, not measured
equality. No delayed follow-up is reported.

The authors report a twenty-five-item multiple-choice pretest. Their method
then says the posttest repeated those items and added six, while giving a
total of twenty-nine and a breakdown of ten inter-level, eight micro-level,
two macro-level and nine other questions. The printed count is internally
inconsistent: twenty-five plus six is not twenty-nine. The following results
are therefore the authors' reported scores and analysis, not a reconstructed
item-denominator calculation.

\begin{table}[htbp]
\centering
\small
\begin{tabular}{p{.39\linewidth}rr}
\hline
2014 condition & Pretest mean & Adjusted posttest mean\\
\hline
Self-explaining, \(n=23\) & 47.0\% & 76.4\%\\
Reading, \(n=19\) & 47.4\% & 68.9\%\\
\hline
\end{tabular}
\caption{The diffusion paper's Table~1 (p.~850). Posttest means are
ANCOVA-adjusted, not raw observed means. The reported comparison is
\(F(1,39)=4.98,p=.03,\eta^2=.086\); the paper does not identify that
effect-size label as partial eta-squared or supply confidence intervals
\cite{MuldnerBurlesonChi2014Emergent}.}
\label{tab:p19-self-explanation-diffusion}
\end{table}

The adjusted comparison favors the implemented explanation package.
The remaining analyses concern only its twenty-three participants.
Explanations were coded as macro-level when they discussed visible patterns,
micro-level when they discussed molecules and interactions, and inter-level
when they explicitly related the two. Verbatim rereading and insufficiently
informative fragments were not explanatory responses. Across all four
categories, including miscellaneous content, 1,087 explanations were counted;
these are coded utterances, not 1,087 independent people or trials assigned
to conditions. Inter-level responses were least common.

The within-group gain correlations were negative for macro explanations,
\(r=-.46,p=.029\), and positive for micro and inter-level explanations,
\(r=.43,p=.04\) and \(r=.44,p=.04\). After separating additional
inferences from paraphrases, the micro and inter-level inference correlations
were \(r=.46,p=.03\) and \(r=.42,p=.049\); macro inferences were not
significantly associated with gains, \(p=.23\), and no paraphrase category
was reliably associated with gains
\cite[pp.~851--852]{MuldnerBurlesonChi2014Emergent}.

These are exploratory process associations, not randomized contrasts among
explanation types. The test had few purely macro-level questions according
to the paper's own description. Prior knowledge, allocation of effort,
available opportunities, and the alignment of explanations with assessed
content are alternatives to a direct effect of producing a particular
type. Selecting the high and low tails of the same gain distribution for
further comparisons does not provide independent confirmation on a new
sample. Multiple correlations and selected-tail analyses need those limits
beside their reported thresholds. In particular, the negative macro
correlation does not prove that explaining visible behavior causes harm.

The protocols also contained incorrect inferences: purposeful molecular
movement, coordinated exchange to maintain equilibrium, and importing a
visible direction of spread into each individual molecule's trajectory.
Some students continued to express these ideas while looking at a simulation
showing random motion. The presence of correct instructional information
does not ensure that an added inference is correct or that a conflict will
be noticed and repaired
\cite[pp.~852--853]{MuldnerBurlesonChi2014Emergent}.

\begin{hbexample}{def:p19-self-explanation-diffusion-worked}{Connect levels without attributing a purpose}
Suppose the instructional material describes a dye becoming evenly visible
while its molecules continue moving randomly. Merely restating the even
color is a macro-level paraphrase. Adding that molecules continue moving
but no visible color gradient remains explicitly relates the two levels.
It addresses the possible mistaken inference from an unchanging visible
pattern to stationary molecules.

The explanation should not add that each molecule knows where dye is scarce
and travels there to balance the mixture. That imports a goal and a
micro-level direction not licensed by the given random-motion account.
A checkable follow-up asks whether a molecule can move back toward the
region where dye was initially concentrated, and whether such motion
contradicts a net spreading pattern. The account permits individual
backward trajectories while describing a population-level pattern.

This is an illustrative reconstruction, not a new measured intervention
or a complete mathematical model of diffusion. It preserves the
micro/macro/inter-level distinction tested in the original task
\cite[pp.~849--853]{MuldnerBurlesonChi2014Emergent}.
\end{hbexample}

\section{Candidate mechanisms and competing explanations}
\label{sec:p19-self-explanation-mechanisms}

The three papers motivate several distinct process claims. \emph{Gap filling}
adds a relation or condition that the text leaves unstated. \emph{Integration}
connects the current line to prior text or general knowledge. \emph{Access}
retrieves knowledge a learner possessed but did not articulate earlier.
\emph{Model repair} changes a previously held relation when new material
conflicts with it. Constructing conditions for applying a procedure is
different again from speeding up a procedure already known.

The physics authors discuss situation-specific rules and qualitative
constraint networks as candidates. Their comparison with explanation-based
AI models points to a difficulty: a learner does not necessarily begin with
the complete correct domain theory that a proof-based generalizer assumes.
This is the authors' interpretation of their learner protocols, not a
qualified independent treatment of those AI systems or proof that one model
uniquely generated the observed behavior
\cite[pp.~177--180]{ChiEtAl1989SelfExplanations}.

The biology authors propose construction, integration with old knowledge,
and continuous piecemeal processing as reasons for the gains. Their protocol
example of a learner revising the route from lungs to body supports the
possibility of local model repair. It does not identify a randomized
mediator. General world knowledge and more extensive across-topic linking
remain possible factors. The authors expressly leave uncertain whether the
same activity repairs concepts whose initial and target representations
belong to incompatible ontological categories; the diffusion experiment is
a later bounded test, not a demonstration that every such difficulty is
resolved
\cite[pp.~470--475]{ChiEtAl1994Eliciting}
\cite[pp.~847--853]{MuldnerBurlesonChi2014Emergent}.

About one quarter of the biology explanations were incorrect. Removing
incorrect explanations did not eliminate the high/low association. That
does not show that errors are beneficial or harmless: students were not
randomly assigned to produce true or false explanations. The proposed
benefit of articulating a mistaken belief is that later information might
make the conflict inspectable. Repair still requires an informative source,
recognition of the conflict and a justified revision. Elaborating a false
premise coherently can instead preserve it.

\begin{hbwarning}{def:p19-self-explanation-causal-boundary}{A prompted package is not a measured mechanism}
A randomized comparison of a prompting package can support a causal
interpretation of that package under the design's assumptions. It does not
randomize how many inferences a student produces, which knowledge the
student brings, or whether a particular inference mediates the outcome.
The 1989 outcome groups, the 1994 explanation-count groups, and the 2014
explanation-type correlations answer different questions. Equal-looking
pretest averages, a non-significant covariate, or a selected subgroup cannot
replace the missing assignment or mediation argument.
\end{hbwarning}

\section{Using and evaluating a self-explanation procedure}
\label{sec:p19-self-explanation-procedure}

\begin{hbmethod}{def:p19-self-explanation-procedure}{Make the missing connection explicit and assess it}
Identify the learning target and the outcome criterion before asking for an
explanation. For a worked step, ask what goal it serves, which premise or
principle permits it, and under what changed condition it would fail. For a
sentence, identify what it adds, how it relates to previous information,
and whether it conflicts with an existing model. Keep a paraphrase, a
confidence report and a substantive additional inference distinct.

Check the added claim against the stipulated task and reliable instructional
information. Preserve uncertainty when the reason is not yet established;
do not reward elaboration solely for length. A useful explanation can be
short and local. An unresolved gap may require background instruction or
feedback rather than more unsupported explanation.

To evaluate the intervention, compare it with a defined alternative,
document assignment, prompts, experimenter attention, time, prior knowledge,
scoring and access to materials. Separate repeated-item performance, new
inferential items, transfer to a specified new task, and delayed performance.
Measure correctness as well as the amount and type of explanation. Report
exclusions, coding reliability and uncertainty; an outcome-conditioned group
is not a treatment group. This is a handbook procedure synthesized from the
bounded designs above, not a newly validated training package.
\end{hbmethod}

The existing explanatory-preference chapter
\cref{sec:p19-explanatory-preferences} concerns selecting a satisfying or
likely explanation from a supplied menu. A preference for one cause rather
than two is not the same measure as learning by generating an inference.
The everyday causal-cognition chapter
\cref{sec:p19-rationality-7-2} separates prediction, diagnosis, structure
learning and intervention; explaining a textbook assertion does not
independently establish the causal structure it asserts. Judgments of
learning in \cref{sec:p19-judgments-of-learning} predict later recall and
distinguish monitoring from control. Detecting a missing reason may trigger
study, but the physics monitoring statements are not calibrated numerical
JOLs and do not inherit the delayed-JOL results.

Analogy in \cref{sec:p19-analogy-and-problem-transfer} adds retrieval,
noticing, mapping and adaptation requirements. A learner may understand a
studied example and still fail to notice when it applies without a hint.
The improvement entry's distinctions among practice, near/far transfer,
retention and generalized efficacy in \cref{def:p19-r153-transfer}
therefore remain necessary. Neither an explanation during a trained task
nor an open-book inferential answer alone demonstrates generalized rational
thinking. A favorable score also does not justify every cost of using a
prompting procedure in another setting.

The 1994 paper varies the age group, domain, material, prompting and outcome
from the 1989 study. The 2014 paper tests a further domain and includes
simulations and specific questions. Their convergence is evidence across
these particular implementations, not a set of exact replications holding
everything else fixed. The held articles do not provide the raw participant
records or confidence intervals needed for a pooled reanalysis. Claims
about additional replications in their introductions are not additional
fully qualified experimental records in this chapter.

For subject/task lookup, a success case is an explicit, correct application
condition that enables adaptation when an example's surface or motion
assumption changes, assessed under the actual new-task conditions. A failure
case is repeating a solution step fluently, replacing a missing reason with
an unsupported causal assertion, and claiming transfer because the same
step can be repeated immediately. A second failure is inferring that a
visible equilibrium means no microscopic motion. The reported studies
support investigation of appropriately targeted constructive explanation;
they do not establish a universal benefit, an optimal prompt schedule,
unrestricted far transfer, or a single necessary learning mechanism.

\index{self-explanation!learning}
\index{learning!explanation versus paraphrase}
\index{worked examples!self-explanation}
\index{transfer!self-explanation boundaries}
